\documentclass[11pt]{article}

% Change "review" to "final" to generate the final (sometimes called camera-ready) version.
% Change to "preprint" to generate a non-anonymous version with page numbers.
\usepackage[preprint]{acl}

% --- IEEE-style numeric citations (override acl.sty's author-year natbib defaults) ---
\setcitestyle{numbers,square,sort&compress}

% Standard package includes
\usepackage{times}
\usepackage{latexsym}

% For proper rendering and hyphenation of words containing Latin characters (including in bib files)
\usepackage[T1]{fontenc}
% For Vietnamese characters
% \usepackage[T5]{fontenc}
% See https://www.latex-project.org/help/documentation/encguide.pdf for other character sets

% This assumes your files are encoded as UTF8
\usepackage[utf8]{inputenc}

% This is not strictly necessary, and may be commented out,
% but it will improve the layout of the manuscript,
% and will typically save some space.
\usepackage{microtype}

% This is also not strictly necessary, and may be commented out.
% However, it will improve the aesthetics of text in
% the typewriter font.
% \usepackage{inconsolata}  % cosmetic typewriter font; requires texlive-fonts-extra

%Including images in your LaTeX document requires adding
%additional package(s)
\usepackage{graphicx}

\usepackage{booktabs}
\usepackage{longtable}
\usepackage{array}
\usepackage{xcolor}
\usepackage{colortbl}
\usepackage{caption}
\usepackage{adjustbox}   % for \begin{adjustbox}
\usepackage{natbib}      % for \citeyear  (or use biblatex)
\usepackage{lscape}      % optional: landscape
\usepackage{multirow}
\usepackage{amssymb}

% If the title and author information does not fit in the area allocated, uncomment the following
%
%\setlength\titlebox{<dim>}
%
% and set <dim> to something 5cm or larger.

\title{Literature Review: Multi-Page Visually Rich Document Understanding}

% Author information can be set in various styles:
% For several authors from the same institution:
% \author{Author 1 \and ... \and Author n \\
%         Address line \\ ... \\ Address line}
% if the names do not fit well on one line use
%         Author 1 \\ {\bf Author 2} \\ ... \\ {\bf Author n} \\
% For authors from different institutions:
% \author{Author 1 \\ Address line \\  ... \\ Address line
%         \And  ... \And
%         Author n \\ Address line \\ ... \\ Address line}
% To start a separate ``row'' of authors use \AND, as in
% \author{Author 1 \\ Address line \\  ... \\ Address line
%         \AND
%         Author 2 \\ Address line \\ ... \\ Address line \And
%         Author 3 \\ Address line \\ ... \\ Address line}

\author{Author: Lewei Xu \\
  The University of Western Australia \\
  Supervisors: Dr Yihao Ding, Dr Daochang Liu \\
  Degree: BACS (Honours) \\}

%\author{
%  \textbf{First Author\textsuperscript{1}},
%  \textbf{Second Author\textsuperscript{1,2}},
%  \textbf{Third T. Author\textsuperscript{1}},
%  \textbf{Fourth Author\textsuperscript{1}},
%\\
%  \textbf{Fifth Author\textsuperscript{1,2}},
%  \textbf{Sixth Author\textsuperscript{1}},
%  \textbf{Seventh Author\textsuperscript{1}},
%  \textbf{Eighth Author \textsuperscript{1,2,3,4}},
%\\
%  \textbf{Ninth Author\textsuperscript{1}},
%  \textbf{Tenth Author\textsuperscript{1}},
%  \textbf{Eleventh E. Author\textsuperscript{1,2,3,4,5}},
%  \textbf{Twelfth Author\textsuperscript{1}},
%\\
%  \textbf{Thirteenth Author\textsuperscript{3}},
%  \textbf{Fourteenth F. Author\textsuperscript{2,4}},
%  \textbf{Fifteenth Author\textsuperscript{1}},
%  \textbf{Sixteenth Author\textsuperscript{1}},
%\\
%  \textbf{Seventeenth S. Author\textsuperscript{4,5}},
%  \textbf{Eighteenth Author\textsuperscript{3,4}},
%  \textbf{Nineteenth N. Author\textsuperscript{2,5}},
%  \textbf{Twentieth Author\textsuperscript{1}}
%\\
%\\
%  \textsuperscript{1}Affiliation 1,
%  \textsuperscript{2}Affiliation 2,
%  \textsuperscript{3}Affiliation 3,
%  \textsuperscript{4}Affiliation 4,
%  \textsuperscript{5}Affiliation 5
%\\
%  \small{
%    \textbf{Correspondence:} \href{mailto:email@domain}{email@domain}
%  }
%}

\begin{document}
\maketitle
\begin{abstract}
Multi-page visually rich document understanding (MP-VRDU) has matured along four architectural families: page-to-document encoders, backbone-centric multimodal language models, retriever-generator pipelines, and agentic pipelines. On top of these, an increasingly elaborate inventory of training-free orchestration techniques has emerged. This review synthesises the literature on both, with disproportionate weight on the retrieval-augmented and agentic strands that drive the current frontier. A single structural assumption recurs across every surveyed system: the retrievable evidence pool is bounded by the input document or, at most, a user-supplied corpus. This review argues that the closed-world assumption is the principal remaining constraint on grounding and out-of-distribution performance in MP-VRDU, and that the agentic substrate already established in the literature is the natural host for an extension that exposes external retrieval as a first-class action alongside intra-document retrieval. The final sections consolidate the open problems specific to the agentic frontier and outline the research direction the project pursues.
\end{abstract}

% =============================================================
\section{Introduction}
% =============================================================

Documents are the principal medium through which organisations, governments, and the scientific community record and exchange information. Contracts, financial reports, scientific papers, technical manuals, regulatory filings, and presentation decks are not only ubiquitous but almost always span multiple pages and interleave dense text with tables, figures, diagrams, and complex layouts~\cite{tito2023hierarchical,borchmann2025arctic}. Manually extracting information from such documents is slow, expensive, and error-prone, motivating a sustained research effort to build automated document visual question answering (DocVQA) systems that can answer natural-language questions over rendered document pages.

Progress on the single-page formulation of this problem has been substantial, with state-of-the-art systems such as Qwen2.5-VL~\cite{bai2025qwen2} and mPLUG-DocOwl1.5~\cite{hu2024mplug1.5} approaching or exceeding human performance on benchmarks such as DocVQA~\cite{mathew2021docvqa}, InfographicVQA~\cite{mathew2022infographicvqa}, and ChartQA~\cite{masry2022chartqa}. The vast majority of real-world documents, however, extend well beyond a single page, and the evidence required to answer a realistic question is typically scattered across several pages or distributed across heterogeneous elements~\cite{tito2023hierarchical,blau2024gram,borchmann2025arctic}. Three structural challenges follow from the move to multi-page inputs. First, the input token count grows linearly (or worse) with page count, rapidly exceeding the context window of even long-context large language models (LLMs) and multimodal LLMs (MLLMs). Second, answering a question often requires resolving cross-page semantic and logical dependencies, including section hierarchy, multi-page tables, and figure-to-text references. Third, explainability becomes considerably harder, as a model must not only produce a correct answer but also surface which pages or passages support it.

Recent surveys cover visually rich document understanding (VRDU) primarily at the single-page level, organising methods by OCR dependency~\cite{ding2025survey,huang2022layoutlmv3}. The multi-page frontier, however, has developed its own architectural taxonomy and an increasingly elaborate inventory of inference-time techniques that are inadequately captured by an OCR-versus-no-OCR axis. This review focuses exclusively on multi-page VRDU and gives proportional weight to the families that now dominate the literature: retriever-generator pipelines and agentic pipelines. The presentation begins with a compact treatment of the four architectural families (\S\ref{sec:architecture}), the modality-specific design challenges that distinguish multi-page from single-page settings (\S\ref{sec:modalities}), and the training (\S\ref{sec:training}) and training-free (\S\ref{sec:trainingfree}) strategies that build on those architectures, with a short summary of the dataset landscape (\S\ref{sec:datasets}). The final two sections then consolidate the research gaps with explicit attention to where agentic pipelines are heading (\S\ref{sec:gaps}) and motivate an external-knowledge-augmented agentic architecture as the specific direction this project pursues (\S\ref{sec:external}).

% =============================================================
\section{Architectural Landscape}
\label{sec:architecture}
% =============================================================

To motivate the specific research direction this project pursues, this review first establishes how the field has organised itself architecturally, since each family carries a different relationship to the evidence pool that ultimately frames the project's contribution. Earlier VRDU work was studied predominantly at the single-page level, and the dominant taxonomy in that setting distinguishes OCR-dependent methods, which consume recognised text and layout boxes, from OCR-free methods, which process page images directly~\cite{ding2025survey,huang2022layoutlmv3}. This OCR axis suffices when the entire document fits in one input but does not resolve the design dimensions introduced by multi-page documents: inputs often exceed the backbone's context window, and the evidence required for a single answer is sparsely distributed across pages and modalities. These pressures force architectural choices about \emph{how} and \emph{where} cross-page evidence is selected and integrated, largely independently of OCR usage. Read through this evidence-management lens, the multi-page literature divides into four recurring families: page-to-document encoding models, backbone-centric MLLM adaptation, retriever-generator pipelines, and agentic pipelines. The four are presented below in roughly the order they entered the field, and the limits of each motivate the next.

\subsection{Page-to-Document Encoding Models}
Page-to-document encoding models process each page through a shared layout-aware encoder such as LayoutLMv3~\cite{huang2022layoutlmv3} or DocFormerv2~\cite{appalaraju2024docformerv2} that fuses OCR text, bounding-box layout, and visual features within the page, then connect per-page representations through cross-page attention~\cite{blau2024gram,borchmann2025arctic}, per-page summary tokens~\cite{tito2023hierarchical}, or recurrent state propagation~\cite{dong2024multi}. By preserving page boundaries and cross-page reading order in the architecture itself, this design provides a strong inductive bias for documents with regular structural templates such as multi-page forms and reports. Every page must nonetheless be jointly encoded in a single forward pass, and the page-to-document compression step discards fine-grained evidence as inputs grow~\cite{blau2024gram,dong2024multi}. Subsequent families address these pressures by replacing the dedicated cross-page encoder with more flexible mechanisms.

\subsection{Backbone-Centric MLLM Adaptation}
Backbone-centric MLLM adaptation removes the dedicated cross-page encoder entirely, instead processing the document inside a pretrained multimodal LLM whose context window absorbs cross-page information directly. Adaptation strategies range from training only lightweight projectors~\cite{tanaka2024instructdoc} or visual-token compression modules~\cite{hu2025mplug2}, to fine-tuning the backbone itself through continued pretraining~\cite{duan2025docopilot}, instruction tuning, or reinforcement learning~\cite{xiong2026docr1}. This design avoids external retrieval or orchestration components and benefits directly from the MLLM's instruction-following and reasoning capabilities. However, the backbone's context length now caps total document size, fine-grained signals depend on the fidelity of token compression~\cite{jia2024leopard}, and document structure such as section hierarchy and figure-caption linkage must be recovered implicitly from pixels~\cite{hannan2025docslm}. These limits motivate moving evidence selection out of the backbone altogether.

\subsection{Retriever-Generator Pipelines}
Retriever-generator pipelines decouple evidence selection from answer generation: a retrieval module returns a compact subset of textual chunks, pages, or graph nodes, and a generator produces an answer conditioned on those units~\cite{cho2024m3docrag,tanaka2025vdocrag}. Because only the selected subset enters the generator, the family lifts the context-length and token-compression limits of backbone-centric models and supports documents that exceed any single backbone's context window. The retrieval and generation stages may be implemented by single modules or distributed across multiple specialised components~\cite{wang2024knowledge,duan2025m2rag,han2025mdocagent}, but the overall control flow remains a fixed pipeline executed once per query, which exposes the retrieved evidence as a partial audit trail. The design assumes that relevant evidence is rankable by query similarity in a single retrieval step, so evidence missed at retrieval is unrecoverable at generation, and questions that require aggregating multiple weakly relevant pages frequently fail~\cite{wu2025molorag}; performance is further bounded by upstream component quality, since chunking, OCR, and modality mismatches propagate directly into the answer. These pressures motivate replacing the fixed pipeline with adaptive, multi-step reasoning.

\subsection{Agentic Pipelines}
Agentic pipelines replace the fixed retrieval pipeline with iterative evidence acquisition: one or more LLM controllers issue successive tool calls for retrieval, navigation, cropping, or OCR, and each subsequent decision is conditioned on the outcomes of earlier observations~\cite{zheng2026docvstar,sun2025docagent}. Roles such as inspection, synthesis, and adjudication may be distributed across heterogeneous backbones~\cite{yu2025mact,wang2025vidorag}, but the defining property is that the trajectory itself adapts to question difficulty rather than following a fixed control flow; multi-component systems with predetermined execution order, even when described as ``agents'', remain within the retriever-generator family. By lifting the single-retrieval-step bottleneck of retriever-generator systems, this design allows compute to scale with question difficulty rather than document length and exposes the trajectory as an inspectable audit trail. However, the variable-length action sequence makes inference cost difficult to bound before execution, controllers remain sensitive to prompt phrasing and may commit early to wrong evidence with silent failure, reproducibility is complicated by stochastic generation, and the system inherits the recall ceiling of every retriever, parser, or OCR module it calls.

Across all four families, the retrievable or attendable evidence pool is bounded by the input document; this shared assumption returns as the central pivot in \S\ref{sec:gaps}. Before reaching that point, the next three sections review how the families represent multimodal evidence (\S\ref{sec:modalities}), how they are trained (\S\ref{sec:training}), and the training-free retrieval and reasoning machinery that has come to define the retriever-generator and agentic frontier (\S\ref{sec:trainingfree}).

% =============================================================
\section{Multi-Page Multimodal Representation}
\label{sec:modalities}
% =============================================================

The architectural choice of \emph{where} cross-page evidence is processed is only half the story; \emph{what} that evidence looks like in each modality drives the rest. Multi-page settings strain text, visual, and structural representation in ways that single-page systems, adequately covered by prior surveys~\cite{ding2025survey}, do not face. The three subsections below identify, for each modality, the multi-page-specific pressure and the small number of design responses that recur across the literature. The recurring response to all three pressures is some form of selective reading rather than uniform encoding, which is the substrate the project's agentic direction inherits.

\subsection{Text Modality}
The pressures specific to multi-page text are twofold: documents often exceed the backbone's context window, and even when they fit, the answer occupies a small fraction of a long input where attention is diluted and key statements are easily lost in the middle~\cite{liu2024lost}. Responses cluster into three patterns. \textbf{Per-page aggregation} encodes each page in detail and combines the results into compact cross-page summaries through attention or memory tokens~\cite{tito2023hierarchical,blau2024gram,borchmann2025arctic}. \textbf{Whole-document compression} keeps every page within a single backbone context by aggressively reducing per-page text tokens through resamplers, layout-aware compressors, or token-reduction modules~\cite{hu2025mplug2,yu2024texthawk2}. \textbf{Retrieve-then-read} abandons full-document encoding entirely, indexing text at chunk or page level and passing only a retrieved subset to the reader; agentic variants extend this into multiple selection-reading rounds~\cite{zhang2024cream,shin2025multidocfusion,sun2025docagent,han2025mdocagent}. Each pattern fails in a characteristic way: per-page aggregation discards token-level detail at the summary step, whole-document compression sacrifices fine-grained reading on individual pages, and retrieve-then-read makes evidence missed at retrieval unrecoverable at the reading stage. The third failure mode is the most consequential for the present project, since it is shared by every retriever-generator and agentic system surveyed in \S\ref{sec:trainingfree}.

\subsection{Visual Modality}
Multi-page visual encoding must remain fine-grained enough to recover charts, tables, and handwritten content that OCR cannot, yet scalable enough to apply across documents spanning tens or hundreds of pages. This tension is absent in the single-page setting, where one page can simply be encoded at the resolution required. The literature manages it through three patterns. \textbf{Aggressive token compression} keeps every page available in context but reduces per-page visual tokens through resamplers or convolutional reducers~\cite{hu2025mplug2,yu2024texthawk2}, preserving coverage at the cost of per-page detail. \textbf{Selective high-resolution reading} restricts expensive encoding to a small subset of pages surfaced by retrieval or controller-driven inspection, sending those high-fidelity tokens into the reader while leaving other pages unread~\cite{zheng2026docvstar,jain2025simpledoc}. \textbf{Retrieval-only visual embeddings} decouple encoding from reading altogether: page images are encoded purely to produce retrieval embeddings, and detailed reading of selected pages is handed to a separate text or vision module~\cite{cho2024m3docrag,tanaka2025vdocrag}. The corresponding failure modes parallel those of the text modality: compression loses fine elements, selective reading inherits the recall ceiling of the selector, and retrieval-only embeddings rank layout-similar pages without confirming that the answer-bearing region is present. They reinforce the same theme that every method here treats the input document as the universe of visual evidence.

\subsection{Layout and Cross-Page Structure}
The structural problem distinctive to multi-page VRDU is \emph{cross-page} structure: section hierarchies that span multiple pages, information referenced or continued across page boundaries, and document-level relationships such as figure-to-text or table-to-caption linkages that no longer fit within a single page's view. Within-page mechanisms inherited from single-page VRDU, such as 2D positional encodings, bounding-box coordinates, and reading-order signals attached to individual tokens~\cite{ding2025survey}, describe only per-page geometry, and per-token boxes do not encode cross-page relationships. End-to-end pretraining of cross-page structure has been attempted~\cite{hu2025mplug2,duan2025docopilot} but is limited by the cost of supervision at this scale. The dominant route is therefore to expose parsed structural information to the model directly, through three patterns. \textbf{Structure-aware chunking} uses a parser to extract section trees and heading hierarchies, then segments documents along semantic boundaries so retrieved units respect section scope rather than sliding-window cuts~\cite{shin2025multidocfusion}. \textbf{Hierarchical retrieval indices} route queries first to coarse units such as sections or chapters and then to fine-grained chunks within the selected branch~\cite{gong2025mhier,zhang2026mldocrag}. \textbf{Graph-based traversal} converts headings, entities, and cross-references into a knowledge graph, allowing controllers to follow typed edges rather than relying on dense similarity alone~\cite{wang2024knowledge}. All three depend on parser quality, so errors in heading detection, table-boundary segmentation, or reading-order extraction propagate directly into every downstream retrieval and navigation step. This dependence is a structural fragility the project inherits, but, as \S\ref{sec:gaps} argues, it is not the limit this project sets out to address.

% =============================================================
\section{Training Strategies}
\label{sec:training}
% =============================================================

How multi-page systems are trained tracks closely with the architectural family they sit in, and the literature has converged on a small set of strategies. This review treats them compactly: the trajectory of interest for this project is the move from monolithic supervised pretraining toward retrieval- and reasoning-targeted objectives, which sets up the training-free orchestration that follows in \S\ref{sec:trainingfree}.

\paragraph{Parameter-Efficient Tuning on Frozen Backbones.}
The most common strategy reuses strong pretrained LLMs or large vision-language models (LVLMs) and updates only a small number of task-specific components such as LoRA~\cite{hu2021lora}, projectors, adapters, layout tokenizers, or task-specific heads, leaving the bulk of pretrained parameters frozen. LayTokenLLM~\cite{zhu2025simple} keeps the LLM backbone frozen and introduces a lightweight layout-tokenisation module together with LoRA-based adaptation; CREAM~\cite{zhang2024cream} extends a similar idea to retrieval-augmented settings; and MoLoRAG~\cite{wu2025molorag} fine-tunes only the retrieval engine over a page graph.

\paragraph{Curriculum and Multi-Stage Training.}
Direct adaptation of single-page models to multi-page objectives tends to underperform on cross-page dependencies, so many works adopt progressive curricula that start from simpler local tasks. DocSLM~\cite{hannan2025docslm} moves from image-OCR alignment through image- and document-level compression to instruction tuning; mPLUG-DocOwl2~\cite{hu2025mplug2} follows a related OCR-free staged strategy, first ensuring that compressed visual tokens preserve text and layout, then pretraining on multi-page parsing.

\paragraph{Self-Supervised Retrieval Training.}
Because retrieval-augmented pipelines are bounded by retriever quality, a growing line of work trains dedicated retrieval components directly. VDocRAG~\cite{tanaka2025vdocrag} introduces self-supervised representation-compression objectives that align dense visual document representations with OCR-derived textual content, and DFVC~\cite{qian2026accurate} trains only a lightweight cross-page fusion adapter on top of a frozen VisRAG or ColQwen backbone. This subline matters here because the visual retrievers it produces are the same components that an external-knowledge agentic system would deploy against a non-document corpus.

\paragraph{Reinforcement Learning for Reasoning.}
Most recent and most directly relevant to the agentic frontier, several systems use reinforcement learning to optimise multi-step reasoning. DocR1~\cite{xiong2026docr1} introduces Evidence Page-Guided GRPO, rewarding the model for first identifying evidence pages and then generating an answer; Doc-V$^{*}$~\cite{zheng2026docvstar} combines supervised trajectory distillation with GRPO over a composite reward; and MACT~\cite{yu2025mact} extends RL to multi-agent settings with mixed per-agent and global rewards. These methods demonstrate that reasoning-quality signals, and not only answer correctness, can be optimised end-to-end, which becomes a prerequisite for training agents whose tools include both intra-document and external retrieval.

% =============================================================
\section{Training-Free Strategies}
\label{sec:trainingfree}
% =============================================================

Training-free orchestration is the locus of the current MP-VRDU frontier and the substrate on which this project builds, so it warrants the most detailed treatment of any section in this review. These strategies improve performance without updating backbone parameters, ranging from offline document-side preparation (parsing into a graph, tree, or outline) to inference-time orchestration through prompting, tool use, retrieval refinement, and multi-agent coordination. They dominate the retriever-generator and agentic families introduced in \S\ref{sec:architecture}, and the project's contribution, an external-retrieval action inside an agentic loop, extends one of the patterns described below.

\subsection{Retrieval and Navigation Techniques}
Training-free retrieval uses off-the-shelf encoders, document structure, or LLMs as routers and re-rankers to select page or chunk evidence. Because retrieval quality bounds every downstream reasoning method, the choice of retrieval signal dominates system behaviour on long documents, and it is the most important locus at which an external evidence source could be introduced.

\paragraph{Sparse Hierarchy-Aware Search.}
The earliest training-free retrieval strategy in this corpus parses the document into an explicit textual structure such as section trees, heading hierarchies, or tree-formatted outlines, and retrieves through term- or structure-anchored matching. Retrieved chunks preserve their parent context, and section-level search and fetch tools allow controllers to navigate the parsed hierarchy directly~\cite{shin2025multidocfusion,sun2025docagent}. More recent variants extend the structure to dual hierarchies covering in-page chunks and cross-page topological summaries, addressing the fragmentation that single-tree designs leave unresolved~\cite{gong2025mhier}. The strength is structural awareness without learned components; the weakness is that sparse signals match on lexical or structural overlap and miss semantic connections and paraphrased queries.

\paragraph{Dense Embedding Search.}
Surface-form sensitivity motivated the shift to dense retrieval, where pages or chunks are encoded into learned representations and queries match on semantic similarity. Most MP-VRDU systems adopt vision-language late-interaction encoders such as ColPali~\cite{faysse2024colpali}, ColBERT~\cite{khattab2020colbert}, and ColQwen~\cite{faysse2025colqwen}, which embed the query as a sequence of contextual token vectors and each page as a bag of patch embeddings, scored by MaxSim aggregation~\cite{wu2025docreact,wu2025molorag,han2025mdocagent,jain2025simpledoc,wang2025vidorag,cho2024m3docrag}. Token-to-patch alignment is particularly important in this setting because a single question token may need to align with a figure caption, table cell, or paragraph buried in a long document. Adaptive selectors further scale the number of retrieved pages with the score distribution rather than using a fixed top-$K$~\cite{li2025avir}. Such methods handle paraphrase and cross-modal queries well but remain structure-blind: sections can be split across embedding chunks, and cross-page references are not represented as first-class objects. Critically for this project, these visual retrievers are well-developed and ready-to-use, which is what makes their application to external corpora a near-immediate engineering possibility.

\paragraph{Joint Multi-Signal Search.}
Several systems combine sparse and dense signals at the retrieval stage rather than committing to a single channel. Typical combinations include dense coarse retrieval filtered by iterative LLM-based re-ranking, page embeddings re-ranked against LLM-generated summaries under persistent working memory, and parallel text and image retrieval tracks whose candidates are aggregated by a downstream agent~\cite{zhang2024cream,jain2025simpledoc,han2025mdocagent}. These combinations improve recall but still score each candidate passage in isolation, so retrieval cannot follow logical links, cross-references, or multi-hop dependencies that span passages.

\paragraph{Graph-Based Traversal.}
Graph approaches address the cross-passage gap by indexing the document as a graph and navigating it through edges. Variants differ in how nodes and edges are defined: knowledge graphs over passages and document-structure nodes traversed by an instruction-tuned LLM~\cite{wang2024knowledge}; page graphs in which VLM-assigned logical relevance is combined with semantic similarity for multi-hop traversal~\cite{wu2025molorag}; and per-query graphs whose chunk nodes are anchored to LVLM-generated query nodes so the structure adapts to the question rather than being fixed at index time~\cite{zhang2026mldocrag}. Graph methods reach references that pure scoring cannot, but incur construction overhead, edge-quality dependence, and traversal cost that scales with hop count. All four retrieval patterns share the same closed-corpus assumption identified in \S\ref{sec:architecture}: every index, embedding, and edge is constructed from the input document itself.

\subsection{Inference-Time Reasoning Strategies}
On top of retrieval, the literature reasons about retrieved evidence through prompting, sampling, and action-selection mechanisms applied at inference. The contrast between the two patterns described below, a fixed evidence buffer versus an adaptive tool loop, is the contrast on which the project's agentic direction depends.

\paragraph{Chain-of-Thought and Self-Improvement.}
Chain-of-thought (CoT) prompting~\cite{wei2022chain} elicits intermediate reasoning steps that integrate retrieved text, tables, and figures across pages before producing a final answer. Reasoning runs over a fixed evidence buffer assembled by a prior retrieval step, so the reasoner cannot acquire new evidence once decoding begins. Extensions add structured prompts over hierarchical retrieval output~\cite{gong2025mhier}, self-refinement and self-verification loops that emit refined queries when evidence is judged insufficient~\cite{jain2025simpledoc}, and document-level self-consistency with uncertainty-based selection~\cite{hannan2025docslm}; trained variants internalise the CoT trajectory through reinforcement learning~\cite{xiong2026docr1}. The shared limit is the closed evidence buffer: questions whose evidence sits outside the retrieved set remain unreachable regardless of how the trace is sampled or verified.

\paragraph{ReAct and Tool-Augmented Reasoning.}
ReAct~\cite{yao2022react} replaces the closed buffer with an action loop, alternating reasoning steps with tool calls whose observations feed the next step. In multi-page VRDU the tools retrieve, fetch, or inspect document pages, letting the reasoner gather evidence during inference. Variants instantiate the loop with InfoNCE-guided sub-query refinement~\cite{wu2025docreact}, outline-driven section-level search and fetch~\cite{sun2025docagent}, intent-conditioned action selection trained on trajectories~\cite{zhang2026docdancer}, and thumbnail-level document maps with working-memory summaries~\cite{zheng2026docvstar}. ReAct is the dominant agentic pattern in the corpus and the natural point at which an external-retrieval action would be inserted: every implementation cited above already exposes an arbitrary tool interface, and the loop itself is agnostic to whether a given tool reads inside or outside the document. The shared costs are tool latency that grows with rounds, prompt sensitivity at each step, and fixed-resolution reading of any fetched page.

\subsection{Multi-Agent Decomposition}
A complementary strand decomposes the work across multiple specialised LLM and VLM agents, such as modality readers, retrievers, judges, and synthesisers, coordinated through a pipeline, controller, or critic loop. Three sub-patterns recur, ordered roughly by increasing flexibility.

\paragraph{Modality-Specialised Pipelines.}
The simplest configuration runs dedicated text and image agents in parallel after retrieval, with a synthesiser fusing their outputs~\cite{han2025mdocagent}, or splits work between a text filter, a visual extractor, and a modal fuser that applies confidence-conditional rules across modalities~\cite{duan2025m2rag}. Parallel pipelines improve modality coverage compared to single-track retrieval, but lack an error-correction stage, so an incorrect output from one branch can dominate aggregation when the other branch lacks counter-evidence.

\paragraph{Verification and Adjudication.}
A second sub-pattern inserts an explicit verifier between proposal and answer. Adjudication-based variants sample multiple candidate answers and route them through an adjudicator that selects the most consistent conclusion~\cite{zhu2025doclens}; actor-reviewer variants pair the answering agent with a reviewer that cross-checks through complementary tools and writes disagreements into a shared reflection memory~\cite{sun2025docagent}. Adjudication catches errors that fusion cannot, but the verifier typically draws on the same backbone family as the actor, so shared blind spots persist.

\paragraph{Adaptive Orchestration.}
The most flexible designs treat the loop itself as adaptive, allocating compute per question and per role. ViDoRAG~\cite{wang2025vidorag} cycles a seeker, an inspector, and an answerer with a Gaussian-mixture model setting per-query $K$ bounds from the similarity distribution; MACT~\cite{yu2025mact} decomposes the loop into planning, execution, judgement, and answer agents with agent-wise adaptive test-time scaling, separating the judge from the corrector to avoid the same agent excusing its own errors. Adaptive orchestration improves the accuracy-latency trade-off but pays in latency that scales with both rounds and agents. Notably, none of the three sub-patterns above ever queries evidence outside the input document, even though their architectures are entirely tool-agnostic. This under-explored direction is taken up in \S\ref{sec:gaps} and \S\ref{sec:external}.

% =============================================================
\section{Datasets and Evaluation}
\label{sec:datasets}
% =============================================================

Dataset and benchmark choices both reflect and reinforce the closed-corpus assumption identified in the preceding sections, which is the principal reason this review treats them as a single concern rather than as a neutral inventory.

Pretraining data reuses general-purpose document archives such as IIT-CDIP~\cite{lewis2006iitcdip} and DocStruct4M~\cite{hu2024mplug1.5} originally designed for single-page tasks; genuinely multi-page corpora have only recently emerged through MP-DocStruct1M, Doc-750K, and PaperPDF~\cite{hu2025mplug2,duan2025docopilot,xie2024wukong}. Fine-tuning has shifted to bespoke synthetic generation per model, with strong LVLMs writing question-answer or trajectory data over scraped PDFs because per-page annotation does not scale to long documents~\cite{tanaka2024instructdoc,jain2025simpledoc,xiong2026docr1,zhang2026docdancer}. Benchmarks have moved from mid-length extractive evaluation on MP-DocVQA, DUDE, and SlideVQA~\cite{tito2023hierarchical,van2023document,tanaka2023slidevqa} to long-context multimodal evaluation on MMLongBench-Doc, LongDocURL, and M-LongDoc~\cite{ma2024mmlongbench,deng2025longdocurl,chia2025m}, with retrieval and page-localisation increasingly scored alongside answer accuracy~\cite{hui2024uda,xiong2026docr1,zhu2025doclens}. MP-DocVQA, MMLongBench-Doc, DUDE, DocVQA, and InfographicVQA dominate cross-model comparison.

Two structural features of this landscape matter for the project's framing. First, the benchmark suite assumes the answer is contained in the input document, which means a model that consults external knowledge is not measured for that capability and may be marked wrong even when its answer is grounded in an authoritative external source. Second, open-domain cross-document retrieval remains rare: only M3DocVQA~\cite{cho2025m3docvqa}, ViDoSeek (introduced by ViDoRAG~\cite{wang2025vidorag}), and OpenDocVQA (introduced by VDocRAG~\cite{tanaka2025vdocrag}) test retrieval over a corpus of documents rather than within a single document, and even those corpora are user-supplied and fixed at evaluation time. Together, these features make the closed-world assumption \emph{measured} as well as \emph{assumed}, and they constrain what a project on external-knowledge augmentation can realistically demonstrate without curating evaluation of its own. This point is taken up in \S\ref{sec:external}.

% =============================================================
\section{Open Problems and the Agentic Frontier}
\label{sec:gaps}
% =============================================================

Sections \ref{sec:architecture}--\ref{sec:datasets} characterised how MP-VRDU systems represent documents, how they are trained, and the training-free machinery that has come to define the retriever-generator and agentic frontier. This section consolidates the resulting open problems, with the synthesis directed toward where agentic pipelines are heading and toward the structural assumption they share with every other family in the field.

\subsection{Recent Progress in Agentic Pipelines}
Agentic systems have produced the clearest qualitative gains over the past two years. By replacing the single-pass retrieve-then-generate control flow of retriever-generator pipelines with iterative evidence acquisition, they raise the recall ceiling that bounds non-iterative RAG: Doc-V$^{*}$~\cite{zheng2026docvstar}, ViDoRAG~\cite{wang2025vidorag}, and DocAgent~\cite{sun2025docagent} all report measurable accuracy gains on long-document benchmarks by issuing multiple retrieval and inspection actions per question. They additionally expose the trajectory itself as an inspectable audit trail of which pages, regions, and modalities the system consulted, partially addressing the explainability gap that motivated MP-VRDU evaluation in the first place. Adaptive compute allocation, instantiated by MACT's per-agent test-time scaling~\cite{yu2025mact}, ViDoRAG's GMM-based top-$K$~\cite{wang2025vidorag}, and Doc-V$^{*}$'s working-memory navigation~\cite{zheng2026docvstar}, allows the same system to spend a few tool calls on easy questions and many on hard ones rather than committing to a fixed pipeline depth. Explicit verification and adjudication, such as DocLens's sampling-adjudication~\cite{zhu2025doclens} and DocAgent's reviewer-with-reflection-memory~\cite{sun2025docagent}, further catch errors that single-branch fusion cannot.

\subsection{Outstanding Problems}
Several open problems nonetheless persist across the surveyed corpus, and most are sharpened rather than resolved by the move to agentic control.

\paragraph{Document-Scale Evidence Aggregation.}
No surveyed architecture has fully decoupled answer quality from document length. Every family imposes some lossy compromise: page-to-document encoders compress per-page summaries, backbone-centric MLLMs dilute attention across long contexts, retriever-generator pipelines drop everything outside the top-$K$, and agentic systems inherit the recall ceiling of each tool they call~\cite{ma2024mmlongbench}. Compression and retrieval choices are typically fixed at index time and independent of the question, so evidence may be discarded long before reasoning begins. Query-conditional and modality-aware compression that adapts to the question's evidential demands remains an open direction.

\paragraph{Retrieval as the Binding Constraint.}
Retrieval performance now defines the upper bound of answer accuracy across the retriever-generator and agentic families~\cite{liu2024lost}. Broadening the retrieved set introduces noise and attention dilution, while narrowing it makes evidence missed at retrieval unrecoverable. Existing approaches suffer from brittle chunking, OCR and parser error propagation, modality mismatch between retriever and generator, and per-passage scoring that ignores cross-reference and logical links. Iterative retrieval~\cite{wu2025molorag,jain2025simpledoc} partially raises this ceiling, but inherits the recall of every retriever and parser it calls.

\paragraph{Cross-Page Structural Reasoning.}
Within-page layout is largely inherited from single-page VRDU through 2D positional embeddings and spatial pretraining~\cite{ding2025survey}, but cross-page structure, including section hierarchy, multi-page tables, and inter-page references, remains the distinctive multi-page problem. Layout-as-objective approaches demand supervision that does not scale; implicit pixel-based recovery depends on backbone capability and high-resolution input; and explicit parsed structure inherits parser fragility. Agentic systems sit on top of these substrates and assume that some structured parse is available, so they amplify rather than escape these limits.

\paragraph{Evaluation Integrity.}
The shift toward LLM-synthesised training data has compromised evaluation integrity at a structural level. MP-DocVQA appears in roughly 15 of 36 surveyed papers' fine-tuning mixes and 18 of their evaluation suites, and DUDE and DocVQA show similar overlap. Of the papers that introduce a synthesised training corpus, a substantial fraction source their documents from existing evaluation benchmarks, producing a circular loop in which a model is trained on a strong LVLM's behaviour over benchmark PDFs and later evaluated against the same PDFs under metrics frequently judged by the same LVLM. Held-out closed-test benchmarks, contamination declarations covering pretraining, fine-tuning, and backbone leakage, and generator-evaluator separation are all standard practice in adjacent fields but not yet in MP-VRDU.

\paragraph{Reproducibility of Stochastic Trajectories.}
Agentic pipelines produce stochastic trajectories whose tool calls and evidence selections vary across runs, complicating reproduction. Per-paper dataset curations and idiosyncratic metric choices prevent direct system-to-system comparison. Trajectory-level evaluation, including standardised retrieval-plus-grounding metrics alongside answer accuracy, is needed but not yet established.

\paragraph{Domain Narrowness and Multilingual Coverage.}
Both training corpora and benchmarks remain overwhelmingly English-centric and concentrated on a narrow slice of domains: industrial filings (MP-DocVQA), financial reports (TAT-DQA~\cite{zhu2022tatdqa}, MultiHiertt~\cite{zhao2022multihiertt}), scientific papers (PaperPDF, Doc-750K, M-LongDoc), and presentation decks (SlideVQA, ViDoSeek). Healthcare, legal-contract reasoning beyond CUAD, government forms, multilingual documents, and handwritten or historical documents are under-represented~\cite{van2023document,ma2024mmlongbench,deng2025longdocurl}. Reported zero-shot generalisation in this setting often reflects in-distribution evaluation, and true unseen-domain performance is undermeasured and likely poor.

\subsection{A Shared Closed-World Assumption}
Across all of the gaps above, a single structural assumption recurs: the retrievable universe is the input document or, at most, a user-supplied corpus. Backbone-centric models compress the document into a context window; retriever-generator pipelines index its pages or chunks; agentic systems issue tool calls against its parsed outline, embeddings, or page graph. In every case, the evidence pool is closed. Even the open-domain benchmarks~\cite{cho2025m3docvqa,wang2025vidorag,tanaka2025vdocrag} are open only over a fixed corpus supplied at evaluation time.

This assumption is the pivot for the remainder of the review. Agentic systems have established the machinery for adaptive evidence acquisition, namely tool-call interfaces, verification, adjudication, working memory, and adaptive compute, but every tool in their inventories points inside the document. A natural extension is to broaden the evidence pool itself, allowing the agent to consult sources outside the document when intra-document evidence is insufficient. Such an extension simultaneously addresses domain narrowness (an unseen domain can be compensated for by querying external knowledge), grounding (external authoritative sources provide attribution beyond the document's own claims), and evidence aggregation (additional retrieval channels relieve the recall ceiling). The following section develops this direction in detail.

% =============================================================
\section{Research Direction for MP-VRDU}
\label{sec:external}
% =============================================================

\subsection{The Closed-World Assumption in MP-VRDU}
Every retrieval system surveyed in \S\ref{sec:trainingfree} indexes content that lives inside the input document. Retrievers such as ColPali~\cite{faysse2024colpali} and the dense visual retrievers built on its lineage~\cite{cho2024m3docrag,tanaka2025vdocrag} operate over page images supplied by the user; sparse and hierarchical indices~\cite{shin2025multidocfusion,gong2025mhier} parse the same document; graph-based traversal~\cite{wang2024knowledge,wu2025molorag,zhang2026mldocrag} constructs its nodes from the document's own content. The handful of open-domain MP-VRDU benchmarks, including M3DocVQA~\cite{cho2025m3docvqa}, ViDoSeek~\cite{wang2025vidorag}, and OpenDocVQA~\cite{tanaka2025vdocrag}, broaden the retrievable unit from a single document to a corpus of documents, but the corpus is still fixed at evaluation time and provided by the user. No surveyed system issues a query against an external knowledge base, a structured authority such as Wikipedia or a domain-specific glossary, or the open web during inference.

This assumption diverges from many practical document question-answering scenarios. A reader of a financial filing may need to expand an unfamiliar accounting acronym; a reader of a regulatory document may need to consult the parent standard the filing references; a reader of a scientific paper may need to verify a numerical claim against the cited source; a reader of a legal contract may need to resolve a jurisdiction-specific term whose definition is not given in the document itself. In each case, the relevant information is not present in the document, so a closed-world system cannot recover the correct answer. Given the increasingly capable language models that underpin modern MP-VRDU agents, the more likely failure mode is confabulation of a plausible answer rather than principled abstention.

\subsection{External-Knowledge RAG in Adjacent Fields}
Adjacent literatures have addressed analogous problems for some time. In text-only question answering, retrieval-augmented generation has been an explicit research programme since at least 2020. The original RAG formulation~\cite{lewis2020rag} and REALM~\cite{guu2020realm} train a language model to condition generation on retrieved Wikipedia passages, demonstrating both factuality gains on long-tail entities and improved generalisation to topics absent from the model's parametric memory. Fusion-in-Decoder~\cite{izacard2021fid} and ATLAS~\cite{izacard2022atlas} extend the paradigm to longer contexts and few-shot settings, and In-Context RALM~\cite{ram2023incontext} shows that retrieval augmentation works with frozen LLMs and arbitrary retrievers, removing the need for specialised pretraining.

A parallel line of work targets the open web directly. WebGPT~\cite{nakano2021webgpt} trains a model to issue search-engine queries and cite the returned pages in its answers, reporting measurable factuality and attribution improvements over a non-retrieval baseline. Internet-augmented dialogue~\cite{komeili2022internet} extends the approach to conversational settings, and Toolformer~\cite{schick2023toolformer} shows that LLMs can be self-trained to invoke external tools, including search APIs, calculators, and translators, from self-supervised data. ReAct~\cite{yao2022react}, the same framework now adopted widely in MP-VRDU agents, was first demonstrated on text question answering by interleaving reasoning with a Wikipedia search API.

Multimodal retrieval-augmented vision-language models have followed a similar trajectory. MuRAG~\cite{chen2022murag} augments a vision-language model with a multimodal external memory of image-text pairs, querying both modalities at inference. REVEAL~\cite{hu2023reveal} retrieves from a heterogeneous corpus of Wikipedia, web images, image-caption pairs, and visual question answering datasets, reporting strong gains on knowledge-intensive visual question answering. These systems indicate that the retrieval-augmentation pattern transfers cleanly into the multimodal setting and yields the same documented benefits in factuality, hallucination reduction, attribution, and robustness on topics absent from training.

\subsection{Obstacles to External Retrieval in MP-VRDU}
The infrastructure required to introduce external retrieval into MP-VRDU is already largely in place. Visual retrievers such as ColPali~\cite{faysse2024colpali}, ColQwen~\cite{faysse2025colqwen}, and VisRAG~\cite{yu2024visrag} are mature and widely adopted inside MP-VRDU pipelines (see \S\ref{sec:trainingfree}), and the same retriever can in principle be applied to an external corpus of rendered web pages, slide-deck mirrors, or PDF archives. Agentic loops with arbitrary tool inventories are similarly established: DocAgent~\cite{sun2025docagent} already orchestrates five distinct retrieval and inspection tools, and adding a \texttt{web\_search} or \texttt{lookup\_entity} action is a small infrastructural extension to its actor-reviewer scaffold. ReAct~\cite{yao2022react}, the dominant agentic pattern in MP-VRDU, was specifically introduced to compose reasoning with arbitrary tool calls including search.

The gap is therefore empirical rather than architectural, and three plausible explanations stand out. First, the benchmark suite that drives the field rewards closed-corpus performance: MMLongBench-Doc~\cite{ma2024mmlongbench}, LongDocURL~\cite{deng2025longdocurl}, MP-DocVQA~\cite{tito2023hierarchical}, and DUDE~\cite{van2023document} all assume the answer is contained in the input document, so external retrieval offers no measured upside on these evaluations and may be penalised when the model is right for an unintended reason. Second, open-domain visual retrieval over the web is genuinely harder than visual retrieval over a fixed corpus, since there is no canonical index, the document distribution is unbounded, and freshness and noise become first-order problems. Third, external retrieval introduces new failure modes such as low-quality results, prompt injection through retrieved content, latency, and dependency on third-party search infrastructure, none of which have yet been studied systematically in this setting.

These obstacles are real but not prohibitive. The text-only RAG community has developed practical responses to each one (retrieval quality estimation, source-attribution metrics, defensive prompting, caching), and the same instruments can be adapted to MP-VRDU. The empirical question is whether the cost is justified by the benefit in the particular setting of multi-page documents, and on what subset of questions external retrieval provides the largest gains.

\subsection{Research Direction}
This project proposes an agentic MP-VRDU system that exposes external retrieval as a first-class action alongside intra-document retrieval, with a routing policy that decides when external evidence is warranted based on intra-document retrieval confidence, terminology absent from the document, or detected domain shift. The system builds on the agentic substrate established by Doc-V$^{*}$~\cite{zheng2026docvstar}, DocAgent~\cite{sun2025docagent}, ViDoRAG~\cite{wang2025vidorag}, and DocDancer~\cite{zhang2026docdancer}, adding a \texttt{web\_search} action and an evidence-reconciliation step that combines intra-document and external evidence into a single grounded answer with explicit attribution. The expected payoffs align directly with the open problems consolidated in \S\ref{sec:gaps}: stronger performance on unseen domains through external knowledge compensation, grounded answers through external attribution metrics, and a principled way to allocate retrieval compute across two complementary evidence sources.

Evaluation should explicitly separate in-distribution and unseen-domain accuracy, report grounding and attribution quality alongside answer correctness, and quantify the new failure modes that external retrieval introduces. Existing benchmarks, including MMLongBench-Doc~\cite{ma2024mmlongbench}, LongDocURL~\cite{deng2025longdocurl}, and the open-domain M3DocVQA~\cite{cho2025m3docvqa} and OpenDocVQA~\cite{tanaka2025vdocrag} suites, support in-distribution evaluation. An additional held-out domain split, drawn from documents whose terminology and conventions are absent from the surveyed training corpora, supports unseen-domain evaluation. A small set of questions deliberately requiring external knowledge, such as unfamiliar acronyms, references to external standards, or claims about entities not described in the document, further isolates the contribution of external retrieval from the contribution of the agentic substrate itself.

The contribution targeted by this project is narrow but specific: the first composition, in MP-VRDU, of intra-document visual retrieval with external (web) retrieval inside an adaptive agentic loop, with grounding and unseen-domain performance as the principal evaluation axes. The literature surveyed above indicates that every component required for this composition already exists; what is missing is the assembly and the evaluation that establishes its empirical value.

% =============================================================
\section{Conclusion}
% =============================================================

Multi-page visually rich document understanding has consolidated around a clear architectural taxonomy and an increasingly elaborate body of training-free orchestration, with the retriever-generator and agentic families dominating the current frontier. The combined machinery of these families supplies adaptive evidence acquisition, modality-aware retrieval, and verification-based reasoning. A single structural assumption nonetheless persists across every family surveyed in this review: the retrievable evidence pool is bounded by the input document or a user-supplied corpus. This closed-world assumption constitutes the principal remaining constraint on grounding and out-of-distribution performance in MP-VRDU. The agentic substrate established in the literature is well-positioned to accommodate an extension that relaxes this constraint, and the project this review motivates pursues that extension, evaluated against the questions, domains, and grounding criteria on which the limitation is most apparent.

% Bibliography style is set inside acl.sty as acl_natbib.bst; we have replaced
% that file with IEEEtran's content so numeric IEEE-style entries are produced.
\bibliography{custom}

\appendix

\end{document}
